\documentclass[../main.tex]{subfiles}
\begin{document}

\chapter{Threads Case Study: PThreads}
\lessoninfo{
  video={P2L3},
  textbook={OSTEP \cite{ostep} ch.~27, 30; Silberschatz et al.\ \cite{silberschatz2018} ch.~4, 7; Kerrisk \cite{kerrisk2010} ch.~29--30},
  keywords={POSIX, POSIX threads, thread creation, thread attributes, joinable thread, detached thread, thread arguments, mutex, trylock, condition variable, safety tips, producer-consumer problem, bounded buffer},
}

Lesson~3 described threads and the mechanisms with which they coordinate in
the abstract notation of Birrell \cite{birrell1989}: \texttt{Fork} and
\texttt{Join}, mutexes and condition variables.  This lesson studies a
concrete interface that provides these mechanisms and that most Unix-like
operating systems support: POSIX threads, or Pthreads.  For each mechanism
it presents the corresponding data types and calls, the details that a real
interface adds to Birrell's model---thread attributes, detached threads,
initialization of synchronization constructs---and rules for using them
safely.  It ends with a complete producer-consumer program.

\section{POSIX threads}
\label{sec:l04-posix}

The \term{Portable Operating System Interface}[POSIX]%
\index{POSIX|see{Portable Operating System Interface}} (POSIX) is a family
of IEEE standards that specifies the interface an operating system provides
to applications.\source{P2L3, lesson preview; POSIX.1 \cite{posix2017};
OSTEP ch.~27.}  It defines the names, arguments and semantics of system
calls and library functions, but not how an operating system implements
them.  Its purpose is portability: a program written against POSIX can be
compiled and run on every operating system that complies with the standard,
not only on the one for which it was first written.

\begin{definition}[POSIX threads]
  The part of POSIX that specifies the interface for creating and managing
  threads, together with the synchronization constructs with which threads
  coordinate---mutexes and condition variables---, is called
  \term{POSIX threads}[POSIXスレッド]\index{Pthreads|see{POSIX threads}}, or
  \emph{Pthreads}.
\end{definition}

Pthreads is part of the base POSIX.1 standard.  Like the rest of POSIX, it
leaves the implementation open: a Pthreads library may, for example, map its
threads onto kernel-level threads according to any of the multithreading
models of Lesson~3.\margin{Linux's Pthreads library is NPTL, which gives
each thread its own kernel task (\texttt{clone(2)}), i.e.\ the one-to-one
model; Lesson~5 examines it.}  The names of its functions and types begin
with \texttt{pthread\_} and those of its constants with
\texttt{PTHREAD\_}; the
data types end in \texttt{\_t} and are opaque, so a program manipulates them
only through the functions of the interface.

The interface closely follows the mechanisms of Lesson~3.
\Cref{tab:l04-birrell} lists the Pthreads counterpart of each of Birrell's
operations; the following sections take them in turn.

\begin{table}[htbp]
  \centering
  \caption{Birrell's mechanisms (Lesson~3) and their Pthreads counterparts.
    \texttt{t}, \texttt{m} and \texttt{c} denote a thread, a mutex and a
    condition variable.}
  \label{tab:l04-birrell}
  \small
  \setlength{\tabcolsep}{5pt}
  \begin{tabular}{@{}lll@{}}
    \toprule
    Operation & Birrell & Pthreads \\
    \midrule
    thread type    & \texttt{Thread}           & \texttt{pthread\_t} \\
    create         & \texttt{Fork(proc, args)}
                   & \texttt{pthread\_create(\&t, \&attr, proc, arg)} \\
    join           & \texttt{Join(t)}          & \texttt{pthread\_join(t, \&status)} \\
    \midrule
    mutex type     & \texttt{Mutex}            & \texttt{pthread\_mutex\_t} \\
    lock           & \texttt{Lock(m) \{}       & \texttt{pthread\_mutex\_lock(\&m)} \\
    unlock         & \texttt{\}}               & \texttt{pthread\_mutex\_unlock(\&m)} \\
    \midrule
    condition type & \texttt{Condition}        & \texttt{pthread\_cond\_t} \\
    wait           & \texttt{Wait(m, c)}       & \texttt{pthread\_cond\_wait(\&c, \&m)} \\
    signal         & \texttt{Signal(c)}        & \texttt{pthread\_cond\_signal(\&c)} \\
    broadcast      & \texttt{Broadcast(c)}     & \texttt{pthread\_cond\_broadcast(\&c)} \\
    \bottomrule
  \end{tabular}
\end{table}

\section{Thread creation}
\label{sec:l04-creation}

In Pthreads a thread is represented by a variable of type
\texttt{pthread\_t}, the counterpart of Birrell's
\texttt{Thread}.\source{P2L3, PThread creation; OSTEP ch.~27; Kerrisk
ch.~29.}  It identifies the thread, and through it the library reaches the
thread's data structure with its identifier, its execution state and the
other information needed to manage it.\caution{Opaque also means that two
identifiers are compared with \texttt{pthread\_equal(t1, t2)} and not with
\texttt{==}: POSIX allows \texttt{pthread\_t} to be a structure.}  Threads
are created and joined with
two functions:
\begin{lstlisting}[language=C, numbers=none]
#include <pthread.h>

int pthread_create(pthread_t *thread,
                   const pthread_attr_t *attr,
                   void *(*start_routine)(void *),
                   void *arg);
int pthread_join(pthread_t thread, void **status);
\end{lstlisting}

\texttt{pthread\_create} corresponds to \texttt{Fork(proc, args)}.  The new
thread starts by executing \texttt{start\_routine(arg)}, and the identifier
of the new thread is stored in \texttt{*thread}.  The start routine takes and
returns a \texttt{void~*}, so that a pointer to data of any type can be
passed to the thread; several arguments are packed into a structure and
passed as a pointer to it.  The second argument, \texttt{attr}, has no
counterpart in Birrell's \texttt{Fork}; it is described below.  The call
returns 0 when the thread has been created and an error number otherwise,
for instance when the system lacks the resources for another thread.

\texttt{pthread\_join} corresponds to \texttt{Join(t)}: it blocks the caller
until \texttt{thread} has terminated.  A thread terminates when its start
routine returns or when it calls \texttt{pthread\_exit(status)}.  The value
it returned, or passed to \texttt{pthread\_exit}, is stored in
\texttt{*status}; a caller that does not need the result passes
\texttt{NULL}.  Like \texttt{pthread\_create}, \texttt{pthread\_join}
returns 0 on success and an error number on failure.\margin{A create-and-join
round trip costs on the order of \SI{10}{\micro\second} on current x86-64
Linux; a thread pool (Lesson~3) pays it once per worker.}

\subsection{Thread attributes}
\label{sec:l04-attr}

\begin{definition}[Thread attribute]
  A property of a thread that is set when the thread is created---such as
  its stack size, whether it can be joined, its scheduling policy and
  priority, or its scope---is a \term{thread attribute}[スレッド属性].  The
  attributes of a new thread are passed to \texttt{pthread\_create} in an
  attribute object of type \texttt{pthread\_attr\_t}.
\end{definition}

\Needspace{8\baselineskip}
Passing \texttt{NULL} instead of an attribute object creates a thread with
default attributes.\margin{Default stack: the \texttt{RLIMIT\_STACK} soft
limit, usually \SI{8}{\mebi\byte} of address space (Lesson~3); the minimum
is \texttt{PTHREAD\_STACK\_MIN}, \SI{16}{\kibi\byte} on x86-64
(\texttt{pthread\_attr\_setstacksize(3)}).}  An attribute object is itself
opaque and is manipulated only through functions:
\begin{lstlisting}[language=C, numbers=none]
int pthread_attr_init(pthread_attr_t *attr);
int pthread_attr_destroy(pthread_attr_t *attr);
/* for each attribute xxx: */
int pthread_attr_setxxx(pthread_attr_t *attr, ...);
int pthread_attr_getxxx(const pthread_attr_t *attr, ...);
\end{lstlisting}
\texttt{pthread\_attr\_init} fills the object with the default values; it
must be called before any other function is applied to the object, and
\texttt{pthread\_attr\_destroy} releases it.  The object is consulted only
while \texttt{pthread\_create} runs, so it can be destroyed, or reused for
further threads, as soon as the thread
exists.\margin{\texttt{SCHED\_OTHER} is the CFS of Lesson~7;
\texttt{SCHED\_FIFO} and \texttt{SCHED\_RR} are real time, priorities
1--99, quantum \SI{100}{\milli\second} (\texttt{sched(7)}).}
\Cref{tab:l04-attributes} lists the attributes that matter in this course.  The scope attribute
selects between the system scope and the process scope of Lesson~3: with
\texttt{PTHREAD\_SCOPE\_SYSTEM} the thread competes for the CPU with all
threads in the system, with \texttt{PTHREAD\_SCOPE\_PROCESS} only with the
threads of its own process.\margin{Linux supports only system scope; a
request for \texttt{PTHREAD\_SCOPE\_PROCESS} fails with
\texttt{ENOTSUP}.}\index{system scope}\index{process scope}

\begin{table}[htbp]
  \centering
  \caption{Thread attributes and the functions that set them.  Each has a
    matching \texttt{get} function.}
  \label{tab:l04-attributes}
  \small
  \begin{tabular}{@{}>{\RaggedRight}p{18mm}l>{\RaggedRight\arraybackslash}p{42mm}@{}}
    \toprule
    Attribute & Set with & Meaning \\
    \midrule
    stack size   & \texttt{pthread\_attr\_setstacksize}
                 & size of the thread's stack \\
    detach state & \texttt{pthread\_attr\_setdetachstate}
                 & joinable (default) or detached (\cref{sec:l04-detached}) \\
    scope        & \texttt{pthread\_attr\_setscope}
                 & system or process scope \\
    inheritance  & \texttt{pthread\_attr\_setinheritsched}
                 & inherit policy and priority from the creator, or take
                   them from the object \\
    scheduling policy & \texttt{pthread\_attr\_setschedpolicy}
                 & e.g.\ \texttt{SCHED\_OTHER}, \texttt{SCHED\_FIFO},
                   \texttt{SCHED\_RR} \\
    priority     & \texttt{pthread\_attr\_setschedparam}
                 & priority within the policy \\
    \bottomrule
  \end{tabular}
\end{table}

\Needspace{10\baselineskip}
\section{Detached threads}
\label{sec:l04-detached}

\begin{definition}[Joinable thread]
  A \term{joinable thread}[ジョイン可能スレッド] is a thread that another
  thread can join: when it terminates, its identifier and exit status are
  kept until a \texttt{pthread\_join} collects them.  A Pthreads thread is
  joinable unless it is created detached or detached later.
\end{definition}

The threads of a process are peers: any of them, not only the thread that
created it, may join a joinable thread.\source{P2L3, detaching PThreads;
Kerrisk ch.~29.}  The usual arrangement is nevertheless that a parent
thread creates children and joins them later, and the joins do more than
collect results.  A joinable thread that has terminated but has not been
joined remains a \emph{zombie} thread: its identifier and exit status are
kept until some thread joins it, and if none ever does, they occupy
resources for as long as the process lives.\margin{The same accounting as for
zombie processes (Lesson~2); a server that never joins its workers
eventually cannot create new ones.}  A main thread must in addition
not return from \texttt{main} before its children have finished, because
returning ends the whole process and with it every thread that is still
running.

\begin{definition}[Detached thread]
  A \term{detached thread}[デタッチスレッド] is a thread that cannot be
  joined.  When it terminates, its resources are released immediately.
\end{definition}

\Needspace{6\baselineskip}
A thread can be detached after its creation, by itself or by another
thread, or it can be created detached by setting the detach state in its
attribute object:\tip{A thread that nobody will join detaches itself with
\texttt{pthread\_detach(pthread\_self())} as its first statement.}
\begin{lstlisting}[language=C, numbers=none]
int pthread_detach(pthread_t thread);

pthread_attr_setdetachstate(&attr, PTHREAD_CREATE_DETACHED);
\end{lstlisting}
A detached thread cannot be made joinable again.  Detaching removes the
need to join: nobody has to collect the thread's status, and its resources
are reclaimed as soon as it terminates.  It does not change the
relation between the threads, which are peers in either case.  A thread
ends without ending the process by calling \texttt{pthread\_exit(status)}:
it terminates only the calling thread, and the process continues until its
last thread has terminated, whether the remaining threads are joinable or
detached (\cref{fig:l04-join-detach}).\caution{Returning from \texttt{main} or
calling \texttt{exit} ends the whole process, detached threads included.}

\begin{figure}[htbp]
  \centering
  % Joinable versus detached threads: lifetimes of the main thread and a child.
% Usage: % Joinable versus detached threads: lifetimes of the main thread and a child.
% Usage: % Joinable versus detached threads: lifetimes of the main thread and a child.
% Usage: \input{figures/l04-join-detach}   (width ~116 mm)
\begin{tikzpicture}[x=1mm, y=1mm, font=\footnotesize\sffamily,
  life/.style={line width=2pt, ln-accent},
  gone/.style={draw=ln-muted, fill=ln-box, minimum height=3mm,
               inner sep=0pt, anchor=west},
  ev/.style={-{Stealth[length=1.6mm]}},
  call/.style={font=\scriptsize\ttfamily, anchor=base, inner sep=1pt},
  lab/.style={font=\scriptsize\sffamily, align=center, text=ln-muted},
  ttl/.style={font=\footnotesize\sffamily\bfseries, anchor=base},
  frame/.style={draw=ln-rule, rounded corners=2pt}]
  % ---------------------------------------------------------- joinable
  \draw[frame] (0,-13) rectangle (55,19);
  \node[ttl] at (27.5,21) {\iflangja{ジョイン可能スレッド}{joinable thread}};
  \node[anchor=west] at (1,10) {main};
  \node[anchor=west] at (1,0) {\iflangja{子}{child}};
  \draw[life] (11,10) -- (53,10);
  \draw[ev] (19,10) -- (19,1.6);
  \node[call] at (19,12.6) {pthread\_create};
  \draw[life] (19,0) -- (31,0);
  \node[gone, minimum width=13mm] (z) at (31,0) {};
  \draw[ev] (44,10) -- (44,1.7);
  \node[call] at (44,12.6) {pthread\_join};
  \node[lab, anchor=north] at (37.5,-2.5)
    {\iflangja{終了したが未ジョイン\\（資源を保持）}{terminated, not joined\\(resources kept)}};
  % ---------------------------------------------------------- detached
  \draw[frame] (61,-13) rectangle (116,19);
  \node[ttl] at (88.5,21) {\iflangja{デタッチスレッド}{detached thread}};
  \node[anchor=west] at (62,10) {main};
  \node[anchor=west] at (62,0) {\iflangja{子}{child}};
  \draw[life] (72,10) -- (101,10);
  \draw[thick] (99.6,8.6) -- (102.4,11.4) (99.6,11.4) -- (102.4,8.6);
  \node[call] at (101,12.6) {pthread\_exit};
  \draw[ev] (78,10) -- (78,1.6);
  \node[call] at (78,12.6) {pthread\_create};
  \draw[life] (78,0) -- (111,0);
  \draw[thick] (111,-2) -- (111,2);
  \node[lab, anchor=north] at (100,-2.5)
    {\iflangja{終了時に資源を解放}{resources released\\on termination}};
\end{tikzpicture}
   (width ~116 mm)
\begin{tikzpicture}[x=1mm, y=1mm, font=\footnotesize\sffamily,
  life/.style={line width=2pt, ln-accent},
  gone/.style={draw=ln-muted, fill=ln-box, minimum height=3mm,
               inner sep=0pt, anchor=west},
  ev/.style={-{Stealth[length=1.6mm]}},
  call/.style={font=\scriptsize\ttfamily, anchor=base, inner sep=1pt},
  lab/.style={font=\scriptsize\sffamily, align=center, text=ln-muted},
  ttl/.style={font=\footnotesize\sffamily\bfseries, anchor=base},
  frame/.style={draw=ln-rule, rounded corners=2pt}]
  % ---------------------------------------------------------- joinable
  \draw[frame] (0,-13) rectangle (55,19);
  \node[ttl] at (27.5,21) {\iflangja{ジョイン可能スレッド}{joinable thread}};
  \node[anchor=west] at (1,10) {main};
  \node[anchor=west] at (1,0) {\iflangja{子}{child}};
  \draw[life] (11,10) -- (53,10);
  \draw[ev] (19,10) -- (19,1.6);
  \node[call] at (19,12.6) {pthread\_create};
  \draw[life] (19,0) -- (31,0);
  \node[gone, minimum width=13mm] (z) at (31,0) {};
  \draw[ev] (44,10) -- (44,1.7);
  \node[call] at (44,12.6) {pthread\_join};
  \node[lab, anchor=north] at (37.5,-2.5)
    {\iflangja{終了したが未ジョイン\\（資源を保持）}{terminated, not joined\\(resources kept)}};
  % ---------------------------------------------------------- detached
  \draw[frame] (61,-13) rectangle (116,19);
  \node[ttl] at (88.5,21) {\iflangja{デタッチスレッド}{detached thread}};
  \node[anchor=west] at (62,10) {main};
  \node[anchor=west] at (62,0) {\iflangja{子}{child}};
  \draw[life] (72,10) -- (101,10);
  \draw[thick] (99.6,8.6) -- (102.4,11.4) (99.6,11.4) -- (102.4,8.6);
  \node[call] at (101,12.6) {pthread\_exit};
  \draw[ev] (78,10) -- (78,1.6);
  \node[call] at (78,12.6) {pthread\_create};
  \draw[life] (78,0) -- (111,0);
  \draw[thick] (111,-2) -- (111,2);
  \node[lab, anchor=north] at (100,-2.5)
    {\iflangja{終了時に資源を解放}{resources released\\on termination}};
\end{tikzpicture}
   (width ~116 mm)
\begin{tikzpicture}[x=1mm, y=1mm, font=\footnotesize\sffamily,
  life/.style={line width=2pt, ln-accent},
  gone/.style={draw=ln-muted, fill=ln-box, minimum height=3mm,
               inner sep=0pt, anchor=west},
  ev/.style={-{Stealth[length=1.6mm]}},
  call/.style={font=\scriptsize\ttfamily, anchor=base, inner sep=1pt},
  lab/.style={font=\scriptsize\sffamily, align=center, text=ln-muted},
  ttl/.style={font=\footnotesize\sffamily\bfseries, anchor=base},
  frame/.style={draw=ln-rule, rounded corners=2pt}]
  % ---------------------------------------------------------- joinable
  \draw[frame] (0,-13) rectangle (55,19);
  \node[ttl] at (27.5,21) {\iflangja{ジョイン可能スレッド}{joinable thread}};
  \node[anchor=west] at (1,10) {main};
  \node[anchor=west] at (1,0) {\iflangja{子}{child}};
  \draw[life] (11,10) -- (53,10);
  \draw[ev] (19,10) -- (19,1.6);
  \node[call] at (19,12.6) {pthread\_create};
  \draw[life] (19,0) -- (31,0);
  \node[gone, minimum width=13mm] (z) at (31,0) {};
  \draw[ev] (44,10) -- (44,1.7);
  \node[call] at (44,12.6) {pthread\_join};
  \node[lab, anchor=north] at (37.5,-2.5)
    {\iflangja{終了したが未ジョイン\\（資源を保持）}{terminated, not joined\\(resources kept)}};
  % ---------------------------------------------------------- detached
  \draw[frame] (61,-13) rectangle (116,19);
  \node[ttl] at (88.5,21) {\iflangja{デタッチスレッド}{detached thread}};
  \node[anchor=west] at (62,10) {main};
  \node[anchor=west] at (62,0) {\iflangja{子}{child}};
  \draw[life] (72,10) -- (101,10);
  \draw[thick] (99.6,8.6) -- (102.4,11.4) (99.6,11.4) -- (102.4,8.6);
  \node[call] at (101,12.6) {pthread\_exit};
  \draw[ev] (78,10) -- (78,1.6);
  \node[call] at (78,12.6) {pthread\_create};
  \draw[life] (78,0) -- (111,0);
  \draw[thick] (111,-2) -- (111,2);
  \node[lab, anchor=north] at (100,-2.5)
    {\iflangja{終了時に資源を解放}{resources released\\on termination}};
\end{tikzpicture}

  \caption{Left: a joinable child that terminates before it is joined keeps
    its resources until \texttt{pthread\_join} collects its status.  Right:
    a detached child releases its resources itself; the main thread ends
    with \texttt{pthread\_exit}, which leaves the process running whether
    the remaining threads are joinable or detached.}
  \label{fig:l04-join-detach}
\end{figure}

\Needspace{9\baselineskip}
\begin{example}
  \label{ex:l04-detached}
  The main thread starts a detached logging thread with system scope and
  exits, leaving the logger to finish its work.  The logger is an ordinary
  start routine:
  \begin{lstlisting}[language=C, numbers=none]
void *logger(void *arg) {            /* thread main */
  printf("logger: flushing records\n");
  pthread_exit(NULL);
}
\end{lstlisting}
  \Needspace{14\baselineskip}
  The main thread sets up the attribute object, creates the logger and
  exits:
  \begin{lstlisting}[language=C, numbers=none]
int main(void) {
  pthread_t tid;
  pthread_attr_t attr;

  pthread_attr_init(&attr);          /* required */
  pthread_attr_setdetachstate(&attr,
                              PTHREAD_CREATE_DETACHED);
  pthread_attr_setscope(&attr, PTHREAD_SCOPE_SYSTEM);
  pthread_create(&tid, &attr, logger, NULL);
  pthread_attr_destroy(&attr);
  pthread_exit(NULL);                /* not: return 0 */
}
\end{lstlisting}
  Had \texttt{main} ended with \texttt{return 0}, the process could
  terminate before the logger prints anything.
\end{example}

\begin{keypoint}
  A Pthreads thread is joinable unless it is detached.  Join every joinable
  thread, or detach threads whose result nobody needs; a main thread that
  must not end the process calls \texttt{pthread\_exit} instead of
  returning.
\end{keypoint}

\section{Compiling Pthreads programs}
\label{sec:l04-compiling}

A program that uses Pthreads needs three things.\source{P2L3, compiling
PThreads; OSTEP ch.~27.}
\begin{enumerate}
  \item \textbf{The header.}  Every source file that uses Pthreads includes
    \texttt{<pthread.h>}, which declares the types, constants and
    functions.
  \item \textbf{The compiler flag.}  The program is compiled and linked
    with \texttt{-pthread} or \texttt{-lpthread}:\margin{On GNU/Linux
    \texttt{-pthread} also defines \texttt{\_REENTRANT}; since glibc~2.34
    the Pthreads functions are part of \texttt{libc} itself.}
\begin{lstlisting}[numbers=none]
gcc -o main main.c -pthread
gcc -o main main.c -lpthread
\end{lstlisting}
    \texttt{-lpthread} only links the Pthreads library, whereas
    \texttt{-pthread} also sets the preprocessor definitions that the
    platform requires for multithreaded code; it is the more portable
    choice.
  \item \textbf{Checked return values.}  Pthreads functions report failure
    through their return value, 0 on success and an error number
    otherwise, rather than through \texttt{errno}.\caution{Pthreads is the
    exception in POSIX: system-call wrappers return $-1$ and set
    \texttt{errno}, so \texttt{perror} after a failed
    \texttt{pthread\_create} reports the wrong error.}  A caller that cannot
    continue reports the number, e.g.\ with \texttt{strerror}, and exits.
    A program that ignores the result of \texttt{pthread\_create} may later
    join a thread that was never created, using an identifier that was
    never set.
\end{enumerate}

\section{Thread creation examples}
\label{sec:l04-examples}

Three small programs show how threads are created and joined and how data
reach a new thread.\source{P2L3, PThread creation examples 1--3; OSTEP
ch.~27.}

\Needspace{23\baselineskip}
\subsection{Creating and joining several threads}

The first program creates a fixed number of threads with default
attributes and waits for all of them:
\begin{lstlisting}[language=C, numbers=none]
#include <stdio.h>
#include <pthread.h>
#define NUM_THREADS 3

void *greet(void *arg) {            /* thread main */
  printf("Hello from a worker\n");
  return NULL;
}

int main(void) {
  pthread_t tid[NUM_THREADS];
  int i;

  for (i = 0; i < NUM_THREADS; i++)  /* create */
    pthread_create(&tid[i], NULL, greet, NULL);
  for (i = 0; i < NUM_THREADS; i++)  /* join */
    pthread_join(tid[i], NULL);
  return 0;
}
\end{lstlisting}
The program has two phases: the first loop creates the threads, the second
joins them.  Each worker prints one line.\margin{The lines do not interleave:
POSIX has each stdio call lock its stream (\texttt{flockfile(3)}), so the
\texttt{printf}s need no mutex.}  The workers share no data, so
they need no synchronization, and because all of them print the same text,
the output does not depend on the order in which they are scheduled.  The
joins keep \texttt{main} from returning, and thereby terminating the
process, before every worker has printed its line.

\Needspace{12\baselineskip}
\subsection{Passing an argument}

Suppose each worker should print its own index.  The start routine reads
the index through its argument:\margin{An alternative passes the index by
value, as
\texttt{(void~*)(intptr\_t)i}, and casts it back: no shared variable, but
the conversion is implementation-defined.}
\begin{lstlisting}[language=C, numbers=none]
void *show_index(void *arg) {       /* thread main */
  int *p = (int *)arg;
  int my_num = *p;
  printf("Thread %d\n", my_num);
  return NULL;
}
\end{lstlisting}
\Needspace{4\baselineskip}
An obvious attempt passes it the address of the loop variable:
\begin{lstlisting}[language=C, numbers=none]
for (i = 0; i < NUM_THREADS; i++)
  pthread_create(&tid[i], NULL, show_index, &i);
\end{lstlisting}
The program does not reliably print $0$, $1$ and $2$.\tip{Race detectors find
this without luck: build with \texttt{-fsanitize=thread}, or run under
\texttt{valgrind --tool=helgrind}.}  Every thread receives
the same address, that of the single variable \texttt{i} in the stack frame
of \texttt{main} (\cref{fig:l04-args}, left), and a thread reads the value
only when it executes \texttt{*p}.  By then \texttt{main} may already have
incremented \texttt{i}.  Because \texttt{i} is visible to all threads that
hold its address, a change made by one thread is seen by all of them, and a
thread that reads \texttt{i} while \texttt{main} modifies it is in a data
race (Lesson~3).\index{race condition}  A worker can therefore print an
index that belongs to another worker, or even a value that \texttt{i} takes
only after the creation loop: $3$ when the loop ends, or a value set by the
join loop, which reuses \texttt{i}.  \Cref{tab:l04-arg-race} shows one
possible execution.

\begin{figure}[htbp]
  \centering
  % Passing &i (one variable shared by all threads) versus &t_num[k]
% (one element per thread) as the argument of pthread_create.
% Usage: % Passing &i (one variable shared by all threads) versus &t_num[k]
% (one element per thread) as the argument of pthread_create.
% Usage: % Passing &i (one variable shared by all threads) versus &t_num[k]
% (one element per thread) as the argument of pthread_create.
% Usage: \input{figures/l04-args}   (width ~116 mm)
\begin{tikzpicture}[x=1mm, y=1mm, font=\footnotesize\sffamily,
  th/.style={draw, circle, fill=white, minimum size=7mm, inner sep=0pt},
  cell/.style={draw, fill=ln-box, minimum width=10mm, minimum height=6mm,
               inner sep=0pt, font=\footnotesize\ttfamily},
  ptr/.style={-{Stealth[length=1.6mm]}, thick, ln-accent},
  lab/.style={font=\scriptsize\sffamily, text=ln-muted, align=center},
  ttl/.style={font=\footnotesize\sffamily\bfseries, anchor=base},
  frame/.style={draw=ln-rule, rounded corners=2pt}]
  % ------------------------------------------------------------ shared &i
  \draw[frame] (0,-9) rectangle (55,29);
  \node[ttl] at (27.5,30.8)
    {\texttt{\&i}: \iflangja{全スレッドで共有}{shared by all threads}};
  \node[th] (a0) at (12,21)   {T0};
  \node[th] (a1) at (27.5,21) {T1};
  \node[th] (a2) at (43,21)   {T2};
  \node[cell] (i) at (27.5,3) {?};
  \node[font=\footnotesize\ttfamily, anchor=east] at (i.west) {i\,};
  \foreach \t in {a0,a1,a2} { \draw[ptr] (\t) -- (i); }
  \node[lab, anchor=north] at (27.5,-0.5)
    {\iflangja{main が変更し続ける}{main keeps modifying it}};
  % ------------------------------------------------------ private &t_num[k]
  \draw[frame] (61,-9) rectangle (116,29);
  \node[ttl] at (88.5,30.8)
    {\texttt{\&t\_num[k]}: \iflangja{スレッドごとに別}{one per thread}};
  \node[th] (b0) at (72,21)   {T0};
  \node[th] (b1) at (88.5,21) {T1};
  \node[th] (b2) at (105,21)  {T2};
  \foreach \k in {0,1,2} {
    \node[cell] (c\k) at (78.5+10*\k,3) {\k};
    \node[font=\scriptsize\ttfamily, text=ln-muted, anchor=north]
      at (78.5+10*\k,-0.3) {[\k]};
  }
  \node[font=\footnotesize\ttfamily, anchor=east] at (c0.west) {t\_num\,};
  \draw[ptr] (b0) -- (c0);
  \draw[ptr] (b1) -- (c1);
  \draw[ptr] (b2) -- (c2);
  \node[lab, anchor=north] at (88.5,-4)
    {\iflangja{作成後は変更されない}{not modified after creation}};
\end{tikzpicture}
   (width ~116 mm)
\begin{tikzpicture}[x=1mm, y=1mm, font=\footnotesize\sffamily,
  th/.style={draw, circle, fill=white, minimum size=7mm, inner sep=0pt},
  cell/.style={draw, fill=ln-box, minimum width=10mm, minimum height=6mm,
               inner sep=0pt, font=\footnotesize\ttfamily},
  ptr/.style={-{Stealth[length=1.6mm]}, thick, ln-accent},
  lab/.style={font=\scriptsize\sffamily, text=ln-muted, align=center},
  ttl/.style={font=\footnotesize\sffamily\bfseries, anchor=base},
  frame/.style={draw=ln-rule, rounded corners=2pt}]
  % ------------------------------------------------------------ shared &i
  \draw[frame] (0,-9) rectangle (55,29);
  \node[ttl] at (27.5,30.8)
    {\texttt{\&i}: \iflangja{全スレッドで共有}{shared by all threads}};
  \node[th] (a0) at (12,21)   {T0};
  \node[th] (a1) at (27.5,21) {T1};
  \node[th] (a2) at (43,21)   {T2};
  \node[cell] (i) at (27.5,3) {?};
  \node[font=\footnotesize\ttfamily, anchor=east] at (i.west) {i\,};
  \foreach \t in {a0,a1,a2} { \draw[ptr] (\t) -- (i); }
  \node[lab, anchor=north] at (27.5,-0.5)
    {\iflangja{main が変更し続ける}{main keeps modifying it}};
  % ------------------------------------------------------ private &t_num[k]
  \draw[frame] (61,-9) rectangle (116,29);
  \node[ttl] at (88.5,30.8)
    {\texttt{\&t\_num[k]}: \iflangja{スレッドごとに別}{one per thread}};
  \node[th] (b0) at (72,21)   {T0};
  \node[th] (b1) at (88.5,21) {T1};
  \node[th] (b2) at (105,21)  {T2};
  \foreach \k in {0,1,2} {
    \node[cell] (c\k) at (78.5+10*\k,3) {\k};
    \node[font=\scriptsize\ttfamily, text=ln-muted, anchor=north]
      at (78.5+10*\k,-0.3) {[\k]};
  }
  \node[font=\footnotesize\ttfamily, anchor=east] at (c0.west) {t\_num\,};
  \draw[ptr] (b0) -- (c0);
  \draw[ptr] (b1) -- (c1);
  \draw[ptr] (b2) -- (c2);
  \node[lab, anchor=north] at (88.5,-4)
    {\iflangja{作成後は変更されない}{not modified after creation}};
\end{tikzpicture}
   (width ~116 mm)
\begin{tikzpicture}[x=1mm, y=1mm, font=\footnotesize\sffamily,
  th/.style={draw, circle, fill=white, minimum size=7mm, inner sep=0pt},
  cell/.style={draw, fill=ln-box, minimum width=10mm, minimum height=6mm,
               inner sep=0pt, font=\footnotesize\ttfamily},
  ptr/.style={-{Stealth[length=1.6mm]}, thick, ln-accent},
  lab/.style={font=\scriptsize\sffamily, text=ln-muted, align=center},
  ttl/.style={font=\footnotesize\sffamily\bfseries, anchor=base},
  frame/.style={draw=ln-rule, rounded corners=2pt}]
  % ------------------------------------------------------------ shared &i
  \draw[frame] (0,-9) rectangle (55,29);
  \node[ttl] at (27.5,30.8)
    {\texttt{\&i}: \iflangja{全スレッドで共有}{shared by all threads}};
  \node[th] (a0) at (12,21)   {T0};
  \node[th] (a1) at (27.5,21) {T1};
  \node[th] (a2) at (43,21)   {T2};
  \node[cell] (i) at (27.5,3) {?};
  \node[font=\footnotesize\ttfamily, anchor=east] at (i.west) {i\,};
  \foreach \t in {a0,a1,a2} { \draw[ptr] (\t) -- (i); }
  \node[lab, anchor=north] at (27.5,-0.5)
    {\iflangja{main が変更し続ける}{main keeps modifying it}};
  % ------------------------------------------------------ private &t_num[k]
  \draw[frame] (61,-9) rectangle (116,29);
  \node[ttl] at (88.5,30.8)
    {\texttt{\&t\_num[k]}: \iflangja{スレッドごとに別}{one per thread}};
  \node[th] (b0) at (72,21)   {T0};
  \node[th] (b1) at (88.5,21) {T1};
  \node[th] (b2) at (105,21)  {T2};
  \foreach \k in {0,1,2} {
    \node[cell] (c\k) at (78.5+10*\k,3) {\k};
    \node[font=\scriptsize\ttfamily, text=ln-muted, anchor=north]
      at (78.5+10*\k,-0.3) {[\k]};
  }
  \node[font=\footnotesize\ttfamily, anchor=east] at (c0.west) {t\_num\,};
  \draw[ptr] (b0) -- (c0);
  \draw[ptr] (b1) -- (c1);
  \draw[ptr] (b2) -- (c2);
  \node[lab, anchor=north] at (88.5,-4)
    {\iflangja{作成後は変更されない}{not modified after creation}};
\end{tikzpicture}

  \caption{Left: all threads receive the address of the loop variable
    \texttt{i}, whose value changes while they start.  Right: each thread
    receives the address of its own array element, which does not change
    after the thread has been created.}
  \label{fig:l04-args}
\end{figure}

\begin{table}[htbp]
  \centering
  \caption{One execution of the program that passes \texttt{\&i} with three
    threads.  T0 reads \texttt{i} after \texttt{main} has incremented it and
    T2 reads it after the join loop has reset it; only T1 prints its own
    index, and the program prints 1, 1 and 0.}
  \label{tab:l04-arg-race}
  \small
  \begin{tabular}{@{}clcl@{}}
    \toprule
    Step & Main thread & \texttt{i} & Workers \\
    \midrule
    1 & create T0 with \texttt{\&i}              & 0 & \\
    2 & \texttt{i++}                             & 1 & \\
    3 & create T1 with \texttt{\&i}              & 1 & \\
    4 &                                          & 1 & T0 reads \texttt{*p}: prints 1 \\
    5 &                                          & 1 & T1 reads \texttt{*p}: prints 1 \\
    6 & \texttt{i++}, create T2 with \texttt{\&i} & 2 & \\
    7 & \texttt{i++}, creation loop ends          & 3 & \\
    8 & join loop: \texttt{i = 0}, join T0        & 0 & \\
    9 &                                          & 0 & T2 reads \texttt{*p}: prints 0 \\
    \bottomrule
  \end{tabular}
\end{table}

\Needspace{10\baselineskip}
\subsection{Passing private data}

The race disappears if each thread receives the address of data that
belongs to it alone and that no other thread modifies:\margin{When the creator
cannot wait, the usual arrangement is one argument structure per thread,
allocated with \texttt{malloc} and freed by the thread itself.}
\begin{lstlisting}[language=C, numbers=none]
int t_num[NUM_THREADS];

for (i = 0; i < NUM_THREADS; i++) {
  t_num[i] = i;
  pthread_create(&tid[i], NULL, show_index, &t_num[i]);
}
\end{lstlisting}
Each worker now reads its own element of \texttt{t\_num}, whose value was
set before the thread was created and never changes afterwards
(\cref{fig:l04-args}, right).  Every index is printed exactly once.  The
\emph{order} of the lines, however, is still not determined: it depends on
how the threads are scheduled, and any permutation of the three lines can
occur.\margin{Three lines admit $3! = 6$ orders, so a run that prints them in
order says nothing about the other five.}  The array lives in the stack frame
of \texttt{main}, which does not
return before it has joined all threads, so the memory remains valid for as
long as the threads use it.\caution{Never pass a thread the address of a
local variable of a function that may return before the thread has read
it.}

\begin{keypoint}
  An argument passed by address is shared, not copied.  Give every thread
  its own data---an array element, or a structure allocated for it---that
  outlives the thread's use of it and that no other thread modifies.
\end{keypoint}

\Needspace{9\baselineskip}
\section{Mutexes}
\label{sec:l04-mutex}

Pthreads mutexes solve the mutual exclusion problem among concurrent
threads, as mutexes do in Lesson~3.\source{P2L3, PThread mutexes; OSTEP
ch.~27--28; Silberschatz et al.\ ch.~7.}\index{mutex}  A mutex has type
\texttt{pthread\_mutex\_t}, and the two basic operations are
\begin{lstlisting}[language=C, numbers=none]
int pthread_mutex_lock(pthread_mutex_t *mutex);
int pthread_mutex_unlock(pthread_mutex_t *mutex);
\end{lstlisting}
The semantics are those of Birrell's mutex: a thread that locks a free
mutex becomes its owner, and a thread that locks a mutex owned by another
thread blocks until the mutex is released.\sidenote{Locking a free mutex costs
one atomic read-modify-write instruction in user space: glibc enters the
kernel, through \texttt{futex(2)}, only when a thread must block or must wake
another.  An uncontended lock therefore costs tens of nanoseconds and a
contended one a system call and a context switch.  Drepper,
\enquote{Futexes Are Tricky}, 2011.}  The difference is syntactic, but
it matters.  Birrell's \texttt{Lock(m)~\{\,\ldots\,\}} releases the mutex
implicitly at the end of the block; in Pthreads the critical section is the
code between the two calls, and the program must call
\texttt{pthread\_mutex\_unlock} explicitly on every path out of it,
including early returns and error paths.

\Needspace{16\baselineskip}
\begin{example}
  \label{ex:l04-safe-insert}
  The safe list insertion of Lesson~3 in C with Pthreads:
  \begin{lstlisting}[language=C, numbers=none]
struct node { int value; struct node *next; };

struct node *head = NULL;          /* shared list */
pthread_mutex_t list_m = PTHREAD_MUTEX_INITIALIZER;

void safe_insert(int value) {
  struct node *e = malloc(sizeof *e);
  e->value = value;

  pthread_mutex_lock(&list_m);
  e->next = head;                  /* critical section */
  head = e;
  pthread_mutex_unlock(&list_m);
}
\end{lstlisting}
  Only the two statements that read and write the shared head pointer are
  in the critical section; allocating and filling the new element involves
  no shared data and is done before the lock is taken, which keeps the
  critical section short.
\end{example}

\Needspace{11\baselineskip}
\subsection{Other mutex operations}

Three further operations complete the basic mutex interface:\margin{The types
are small---\SI{40}{\byte} for \texttt{pthread\_mutex\_t} and
\SI{48}{\byte} for \texttt{pthread\_cond\_t} on x86-64 glibc---so a lock
can sit inside each structure it protects.}
\begin{lstlisting}[language=C, numbers=none]
int pthread_mutex_init(pthread_mutex_t *mutex,
                       const pthread_mutexattr_t *attr);
int pthread_mutex_trylock(pthread_mutex_t *mutex);
int pthread_mutex_destroy(pthread_mutex_t *mutex);
\end{lstlisting}
\begin{description}
  \item[\texttt{pthread\_mutex\_init}] initializes a mutex before its first
    use.  Its attribute object, of type \texttt{pthread\_mutexattr\_t},
    specifies the behaviour of the mutex, for instance whether it may be
    shared among processes (Lesson~9); \texttt{NULL} selects the
    defaults.\margin{The type is an attribute too: the default deadlocks when
    its owner locks it again (Lesson~3), \texttt{PTHREAD\_MUTEX\_ERRORCHECK}
    returns \texttt{EDEADLK} instead, \texttt{PTHREAD\_MUTEX\_RECURSIVE}
    counts.}
    A mutex with default attributes can also be initialized statically, as
    in \cref{ex:l04-safe-insert}, with \texttt{PTHREAD\_MUTEX\_INITIALIZER}.
  \item[\texttt{pthread\_mutex\_trylock}] locks the mutex if it is free and
    returns 0.  If the mutex is locked, it does not block but returns
    immediately with the error \texttt{EBUSY}.  The calling thread can then
    do something else and try again later.
  \item[\texttt{pthread\_mutex\_destroy}] frees the data structures of a
    mutex that is no longer needed.  The mutex must not be locked.
\end{description}

\Needspace{8\baselineskip}
\begin{example}
  A thread updates optional statistics only when the mutex happens to be
  free, rather than delaying its main work:\margin{Calling \texttt{trylock} in a
  loop turns it into a spin lock, which burns CPU while another thread holds
  the mutex; Lesson~10 compares spinning with blocking.}
  \begin{lstlisting}[language=C, numbers=none]
if (pthread_mutex_trylock(&stats_m) == 0) {
  stats.requests += local_count;
  pthread_mutex_unlock(&stats_m);
  local_count = 0;
}  /* else: keep counting locally, retry later */
\end{lstlisting}
\end{example}

\subsection{Mutex safety tips}

The most common errors with mutexes are avoided by a few rules.
\begin{itemize}
  \item \textbf{One mutex per piece of shared data.}  All accesses to a
    shared variable must go through the same mutex.  If different code
    paths protect the same data with different mutexes, they exclude
    nothing from each other.
  \item \textbf{A mutex must be visible to all threads that use it.}  It is
    therefore declared at global scope, or in a structure that all these
    threads reach.  A mutex declared as a local variable of the thread
    function is a different mutex in every thread.
  \item \textbf{Lock mutexes in a global order.}  If every thread acquires
    mutexes in the same order, the deadlocks of Lesson~3 cannot occur.%
    \index{lock ordering}\tip{Number the mutexes and always lock in
    increasing order; a thread that then needs a lower-numbered one releases
    what it holds and starts again.}
  \item \textbf{Always unlock, and unlock the correct mutex.}  A forgotten
    unlock blocks every other thread that needs the mutex; unlocking the
    wrong mutex leaves one critical section unprotected and another one
    locked.
\end{itemize}

\section{Condition variables}
\label{sec:l04-condvar}

Pthreads condition variables, of type \texttt{pthread\_cond\_t}, provide
Birrell's three operations.\source{P2L3, PThread condition variables;
OSTEP ch.~27, 30.}\index{condition variable}
\begin{lstlisting}[language=C, numbers=none]
int pthread_cond_wait(pthread_cond_t *cond,
                      pthread_mutex_t *mutex);
int pthread_cond_signal(pthread_cond_t *cond);
int pthread_cond_broadcast(pthread_cond_t *cond);
\end{lstlisting}
The semantics are those of Lesson~3.  \texttt{pthread\_cond\_wait} is called
with \texttt{mutex} held; it atomically releases the mutex and places the
calling thread on the wait queue of \texttt{cond}, and when the thread is
woken up, it re-acquires the mutex before it returns.
\texttt{pthread\_cond\_signal} wakes up at least one of the waiting
threads, \texttt{pthread\_cond\_broadcast} all of them; both have no effect when no
thread is waiting.\caution{The argument order is the reverse of Birrell's
\texttt{Wait(m, c)}: \texttt{pthread\_cond\_wait} takes the condition
variable first.}

POSIX allows \texttt{pthread\_cond\_wait} to return even though no thread
has signalled, and calls such a return a spurious wake-up
too\index{spurious wake-up}; it is a different case from the one examined in
Lesson~3, where a thread is woken by a notification that it cannot use.
Both leave the woken thread unable to proceed, and both are handled by the
same rule: a woken thread tests its predicate again, so the \texttt{while}
loop of Lesson~3 is not merely good practice but required:\sidenote{The
rationale of POSIX.1 \cite{posix2017} explains the licence: requiring that a
thread never return from a wait in vain would make condition variables slower
on a multiprocessor, and an implementation may legitimately implement
\texttt{pthread\_cond\_signal} as a broadcast.  As the predicate has to be
tested anyway, the loop costs nothing and makes a program tolerant of extra
notifications.}
\begin{lstlisting}[language=C, numbers=none]
pthread_mutex_lock(&m);
while (!predicate)
  pthread_cond_wait(&c, &m);
/* predicate holds and m is held */
update_shared_state();
pthread_mutex_unlock(&m);
\end{lstlisting}

\Needspace{6\baselineskip}
Like mutexes, condition variables are initialized and destroyed
explicitly:
\begin{lstlisting}[language=C, numbers=none]
int pthread_cond_init(pthread_cond_t *cond,
                      const pthread_condattr_t *attr);
int pthread_cond_destroy(pthread_cond_t *cond);
\end{lstlisting}
The attribute object of type \texttt{pthread\_condattr\_t} specifies the
behaviour of the condition variable, for instance whether it is shared among
processes; \texttt{NULL} selects the defaults, which the static initializer
\texttt{PTHREAD\_COND\_INITIALIZER} also provides.\margin{Sharing one between
processes requires \texttt{PTHREAD\_PROCESS\_SHARED} on both the condition
variable and its mutex, and both must lie in shared memory (Lesson~9).}

\subsection{Condition variable safety tips}

\begin{itemize}
  \item \textbf{Do not forget to notify waiting threads.}  Whenever a thread
    changes the state that a predicate depends on, it must signal or
    broadcast the condition variable on which the threads that wait for
    this predicate are blocked.
  \item \textbf{When in doubt, broadcast.}  Broadcasting where a signal
    would do is still correct, because every woken thread re-tests its
    predicate and waits again if it cannot proceed; it costs performance,
    since these threads are woken up and scheduled for
    nothing.\margin{With $n$ waiters a broadcast costs $n$ wake-ups and $n$
    mutex acquisitions, $n-1$ of them wasted---the thundering herd of
    Lesson~3.}
  \item \textbf{A mutex is not needed to signal or broadcast.}  When the
    decision to notify does not depend on shared state that can only be read
    under the mutex, the notification can be issued after the mutex has been
    released: the woken threads then do not have to wait for a mutex that
    the signalling thread still holds, which is the cause of spurious
    wake-ups examined in Lesson~3.  Returns from
    \texttt{pthread\_cond\_wait} without any notification remain possible,
    so this does not make the \texttt{while} loop unnecessary.
\end{itemize}

\begin{keypoint}
  Pthreads mutexes and condition variables have Birrell's semantics; what
  they add is explicit unlocking, explicit initialization and destruction,
  and attribute objects.  Wait in a \texttt{while} loop, notify whenever
  the predicate may have changed, and issue the notification after
  unlocking where the decision to notify allows it.
\end{keypoint}

\section{A producer-consumer example}
\label{sec:l04-prodcons}

The producer-consumer problem of Lesson~3 can now be solved with a complete
Pthreads program.\source{P2L3, producer and consumer example using
PThreads; OSTEP ch.~30.}\index{producer-consumer problem}  A producer thread
inserts values into a shared \emph{bounded buffer}, an array of
\texttt{BUF\_SIZE} slots used as a circular queue, and a consumer thread
removes them in the order in which they were inserted.  Four variables
describe the buffer: the array \texttt{buffer}, the index \texttt{add} of the
next slot to fill, the index \texttt{rem} of the next slot to empty, and the
number \texttt{num} of values in the buffer.\margin{\texttt{num} can be
dropped by leaving one slot empty: the buffer is then empty when
\texttt{add == rem} and full when \texttt{(add + 1) \% BUF\_SIZE == rem}.}
Both indices advance modulo
\texttt{BUF\_SIZE}, so the queue wraps around the end of the array
(\cref{fig:l04-bounded-buffer}).

One mutex \texttt{m} protects these variables; the producer waits on
\texttt{c\_prod} while the buffer is full, the consumer on \texttt{c\_cons}
while it is empty.  All shared state is global, as the mutex safety tips
require, and initialized statically:
\begin{lstlisting}[language=C, numbers=none]
#include <stdio.h>
#include <stdlib.h>
#include <pthread.h>

#define BUF_SIZE  4     /* slots in the shared buffer */
#define NUM_ITEMS 10    /* values to produce          */

int buffer[BUF_SIZE];   /* shared circular buffer     */
int add = 0;            /* next slot to fill          */
int rem = 0;            /* next slot to empty         */
int num = 0;            /* values in the buffer       */

pthread_mutex_t m = PTHREAD_MUTEX_INITIALIZER;
pthread_cond_t c_cons = PTHREAD_COND_INITIALIZER;
pthread_cond_t c_prod = PTHREAD_COND_INITIALIZER;

void *producer(void *param);
void *consumer(void *param);
\end{lstlisting}

\begin{figure}[htbp]
  \centering
  % Producer-consumer with a bounded circular buffer of 4 slots,
% one mutex and two condition variables.
% Usage: % Producer-consumer with a bounded circular buffer of 4 slots,
% one mutex and two condition variables.
% Usage: % Producer-consumer with a bounded circular buffer of 4 slots,
% one mutex and two condition variables.
% Usage: \input{figures/l04-bounded-buffer}   (width ~116 mm)
\begin{tikzpicture}[x=1mm, y=1mm, font=\footnotesize\sffamily,
  thr/.style={draw, rounded corners=2pt, fill=white, thick,
              minimum width=21mm, minimum height=8mm, inner sep=1pt},
  cell/.style={draw, fill=white, minimum width=9mm, minimum height=7mm,
               inner sep=0pt, font=\footnotesize\ttfamily},
  full/.style={cell, fill=ln-box},
  >={Stealth[length=1.6mm]},
  tt/.style={font=\scriptsize\ttfamily, inner sep=1pt},
  lab/.style={font=\scriptsize\sffamily, align=center},
  note/.style={font=\scriptsize\sffamily, align=flush left,
               text width=30mm, inner sep=0pt}]
  % ------------------------------------------------ region protected by m
  \draw[draw=ln-accent, dashed, rounded corners=3pt] (34,-15) rectangle (82,14);
  \node[lab, anchor=south, text=ln-accent] at (58,14.5)
    {\iflangja{\texttt{m} が保護: \texttt{buffer}, \texttt{add}, \texttt{rem}, \texttt{num}}%
              {protected by \texttt{m}: \texttt{buffer}, \texttt{add}, \texttt{rem}, \texttt{num}}};
  % ------------------------------------------------------------- buffer
  \node[cell] (s0) at (44.5,0) {};
  \node[full] (s1) at (53.5,0) {7};
  \node[full] (s2) at (62.5,0) {8};
  \node[cell] (s3) at (71.5,0) {};
  \foreach \k in {0,...,3} {
    \node[tt, text=ln-muted, anchor=north] at (44.5+9*\k,-3.8) {\k};
  }
  \draw[->] (71.5,9) node[tt, anchor=south] {add} -- (s3.north);
  \draw[->] (53.5,9) node[tt, anchor=south] {rem} -- (s1.north);
  \draw[->, ln-muted] (s3.south east) .. controls (80,-11) and (36,-11)
    .. (s0.south west);
  \node[tt, anchor=north] at (58,-9.8) {num = 2};
  % ------------------------------------------------------------ threads
  \node[thr] (p) at (13,0)  {\iflangja{生産者}{producer}};
  \node[thr] (c) at (103,0) {\iflangja{消費者}{consumer}};
  \draw[->, thick] (p.east) -- node[lab, above] {\iflangja{挿入}{insert}} (34,0);
  \draw[->, thick] (82,0) -- node[lab, above] {\iflangja{取り出し}{remove}} (c.west);
  \node[note, anchor=north west] at (0,-6)
    {\iflangja{満杯の間 \texttt{c\_prod} で待つ．挿入後に \texttt{c\_cons} を signal}%
              {waits on \texttt{c\_prod} while the buffer is full; signals \texttt{c\_cons} after inserting}};
  \node[note, anchor=north east] at (116,-6)
    {\iflangja{空の間 \texttt{c\_cons} で待つ．取り出し後に \texttt{c\_prod} を signal}%
              {waits on \texttt{c\_cons} while the buffer is empty; signals \texttt{c\_prod} after removing}};
\end{tikzpicture}
   (width ~116 mm)
\begin{tikzpicture}[x=1mm, y=1mm, font=\footnotesize\sffamily,
  thr/.style={draw, rounded corners=2pt, fill=white, thick,
              minimum width=21mm, minimum height=8mm, inner sep=1pt},
  cell/.style={draw, fill=white, minimum width=9mm, minimum height=7mm,
               inner sep=0pt, font=\footnotesize\ttfamily},
  full/.style={cell, fill=ln-box},
  >={Stealth[length=1.6mm]},
  tt/.style={font=\scriptsize\ttfamily, inner sep=1pt},
  lab/.style={font=\scriptsize\sffamily, align=center},
  note/.style={font=\scriptsize\sffamily, align=flush left,
               text width=30mm, inner sep=0pt}]
  % ------------------------------------------------ region protected by m
  \draw[draw=ln-accent, dashed, rounded corners=3pt] (34,-15) rectangle (82,14);
  \node[lab, anchor=south, text=ln-accent] at (58,14.5)
    {\iflangja{\texttt{m} が保護: \texttt{buffer}, \texttt{add}, \texttt{rem}, \texttt{num}}%
              {protected by \texttt{m}: \texttt{buffer}, \texttt{add}, \texttt{rem}, \texttt{num}}};
  % ------------------------------------------------------------- buffer
  \node[cell] (s0) at (44.5,0) {};
  \node[full] (s1) at (53.5,0) {7};
  \node[full] (s2) at (62.5,0) {8};
  \node[cell] (s3) at (71.5,0) {};
  \foreach \k in {0,...,3} {
    \node[tt, text=ln-muted, anchor=north] at (44.5+9*\k,-3.8) {\k};
  }
  \draw[->] (71.5,9) node[tt, anchor=south] {add} -- (s3.north);
  \draw[->] (53.5,9) node[tt, anchor=south] {rem} -- (s1.north);
  \draw[->, ln-muted] (s3.south east) .. controls (80,-11) and (36,-11)
    .. (s0.south west);
  \node[tt, anchor=north] at (58,-9.8) {num = 2};
  % ------------------------------------------------------------ threads
  \node[thr] (p) at (13,0)  {\iflangja{生産者}{producer}};
  \node[thr] (c) at (103,0) {\iflangja{消費者}{consumer}};
  \draw[->, thick] (p.east) -- node[lab, above] {\iflangja{挿入}{insert}} (34,0);
  \draw[->, thick] (82,0) -- node[lab, above] {\iflangja{取り出し}{remove}} (c.west);
  \node[note, anchor=north west] at (0,-6)
    {\iflangja{満杯の間 \texttt{c\_prod} で待つ．挿入後に \texttt{c\_cons} を signal}%
              {waits on \texttt{c\_prod} while the buffer is full; signals \texttt{c\_cons} after inserting}};
  \node[note, anchor=north east] at (116,-6)
    {\iflangja{空の間 \texttt{c\_cons} で待つ．取り出し後に \texttt{c\_prod} を signal}%
              {waits on \texttt{c\_cons} while the buffer is empty; signals \texttt{c\_prod} after removing}};
\end{tikzpicture}
   (width ~116 mm)
\begin{tikzpicture}[x=1mm, y=1mm, font=\footnotesize\sffamily,
  thr/.style={draw, rounded corners=2pt, fill=white, thick,
              minimum width=21mm, minimum height=8mm, inner sep=1pt},
  cell/.style={draw, fill=white, minimum width=9mm, minimum height=7mm,
               inner sep=0pt, font=\footnotesize\ttfamily},
  full/.style={cell, fill=ln-box},
  >={Stealth[length=1.6mm]},
  tt/.style={font=\scriptsize\ttfamily, inner sep=1pt},
  lab/.style={font=\scriptsize\sffamily, align=center},
  note/.style={font=\scriptsize\sffamily, align=flush left,
               text width=30mm, inner sep=0pt}]
  % ------------------------------------------------ region protected by m
  \draw[draw=ln-accent, dashed, rounded corners=3pt] (34,-15) rectangle (82,14);
  \node[lab, anchor=south, text=ln-accent] at (58,14.5)
    {\iflangja{\texttt{m} が保護: \texttt{buffer}, \texttt{add}, \texttt{rem}, \texttt{num}}%
              {protected by \texttt{m}: \texttt{buffer}, \texttt{add}, \texttt{rem}, \texttt{num}}};
  % ------------------------------------------------------------- buffer
  \node[cell] (s0) at (44.5,0) {};
  \node[full] (s1) at (53.5,0) {7};
  \node[full] (s2) at (62.5,0) {8};
  \node[cell] (s3) at (71.5,0) {};
  \foreach \k in {0,...,3} {
    \node[tt, text=ln-muted, anchor=north] at (44.5+9*\k,-3.8) {\k};
  }
  \draw[->] (71.5,9) node[tt, anchor=south] {add} -- (s3.north);
  \draw[->] (53.5,9) node[tt, anchor=south] {rem} -- (s1.north);
  \draw[->, ln-muted] (s3.south east) .. controls (80,-11) and (36,-11)
    .. (s0.south west);
  \node[tt, anchor=north] at (58,-9.8) {num = 2};
  % ------------------------------------------------------------ threads
  \node[thr] (p) at (13,0)  {\iflangja{生産者}{producer}};
  \node[thr] (c) at (103,0) {\iflangja{消費者}{consumer}};
  \draw[->, thick] (p.east) -- node[lab, above] {\iflangja{挿入}{insert}} (34,0);
  \draw[->, thick] (82,0) -- node[lab, above] {\iflangja{取り出し}{remove}} (c.west);
  \node[note, anchor=north west] at (0,-6)
    {\iflangja{満杯の間 \texttt{c\_prod} で待つ．挿入後に \texttt{c\_cons} を signal}%
              {waits on \texttt{c\_prod} while the buffer is full; signals \texttt{c\_cons} after inserting}};
  \node[note, anchor=north east] at (116,-6)
    {\iflangja{空の間 \texttt{c\_cons} で待つ．取り出し後に \texttt{c\_prod} を signal}%
              {waits on \texttt{c\_cons} while the buffer is empty; signals \texttt{c\_prod} after removing}};
\end{tikzpicture}

  \caption{The producer-consumer program with $\texttt{BUF\_SIZE}=4$:
    one mutex protects all buffer state, and each thread waits on its own
    condition variable.}
  \label{fig:l04-bounded-buffer}
\end{figure}

\Needspace{14\baselineskip}
The main thread creates and joins the two threads:
\begin{lstlisting}[language=C, numbers=none]
int main(void) {
  pthread_t tid1, tid2;

  if (pthread_create(&tid1, NULL, producer, NULL) != 0 ||
      pthread_create(&tid2, NULL, consumer, NULL) != 0) {
    fprintf(stderr, "cannot create threads\n");
    exit(1);
  }
  pthread_join(tid1, NULL);   /* wait for producer */
  pthread_join(tid2, NULL);   /* wait for consumer */
  printf("parent done\n");
  return 0;
}
\end{lstlisting}

The producer inserts the values
$1,\dots,\texttt{NUM\_ITEMS}$:\tip{With \texttt{BUF\_SIZE} a power of two,
\texttt{(add + 1) \% BUF\_SIZE} is \texttt{(add + 1) \& (BUF\_SIZE - 1)}---one
instruction instead of a division.}
\begin{lstlisting}[language=C, numbers=none]
void *producer(void *param) {
  int i;
  for (i = 1; i <= NUM_ITEMS; i++) {
    pthread_mutex_lock(&m);
    while (num == BUF_SIZE)         /* full: wait  */
      pthread_cond_wait(&c_prod, &m);
    buffer[add] = i;                /* not full    */
    add = (add + 1) % BUF_SIZE;
    num++;
    pthread_mutex_unlock(&m);
    pthread_cond_signal(&c_cons);   /* not empty   */
    printf("produced %d\n", i);
  }
  return NULL;
}
\end{lstlisting}

The consumer is symmetric:\margin{Each thread has its own \texttt{i} and
\texttt{n} on its own stack (Lesson~3); \texttt{buffer}, \texttt{add},
\texttt{rem} and \texttt{num} are the only shared state.}
\begin{lstlisting}[language=C, numbers=none]
void *consumer(void *param) {
  int i, n;
  for (n = 0; n < NUM_ITEMS; n++) {
    pthread_mutex_lock(&m);
    while (num == 0)                /* empty: wait */
      pthread_cond_wait(&c_cons, &m);
    i = buffer[rem];                /* not empty   */
    rem = (rem + 1) % BUF_SIZE;
    num--;
    pthread_mutex_unlock(&m);
    pthread_cond_signal(&c_prod);   /* not full    */
    printf("consumed %d\n", i);
  }
  return NULL;
}
\end{lstlisting}

The program applies the rules of this lesson and of Lesson~3.
\begin{itemize}
  \item Each thread tests its predicate in a \texttt{while} loop while
    holding \texttt{m}.  When the wait returns, the thread holds \texttt{m}
    again, and it re-tests the predicate before it touches the buffer.
  \item Each thread signals after unlocking, so the woken thread does not
    have to wait for \texttt{m}.\margin{The \texttt{printf} is outside the
    critical section for the same reason: writing to a terminal can block,
    and a thread that blocks while holding \texttt{m} stops the other one.}
    A signal when nobody waits has no effect;
    this is harmless, because a thread tests the predicate before it waits
    and therefore never waits for a change that has already happened.
  \item Because producers wait only on \texttt{c\_prod} and consumers only
    on \texttt{c\_cons}, a notification always reaches a thread of the kind
    that can use it.\caution{With a single condition variable for both, a
    \texttt{signal} could wake a producer when only a consumer can proceed;
    such a design needs \texttt{broadcast}.}
  \item The consumer stops after \texttt{NUM\_ITEMS} values.  A consumer
    that loops forever would never terminate, and the second
    \texttt{pthread\_join} in \texttt{main} would never return.
\end{itemize}

\Needspace{6\baselineskip}
\begin{example}
  \Cref{tab:l04-prodcons-trace} traces one execution of the program in
  which the producer runs ahead and fills the buffer before the consumer
  starts.  Removed values are shown as \texttt{-}.
\end{example}

\begin{table}[htbp]
  \centering
  \caption{A trace of the producer (P) and consumer (C) with
    $\texttt{BUF\_SIZE}=4$.  The buffer column lists slots 0--3; the values
    of \texttt{add}, \texttt{rem} and \texttt{num} are those after each
    step.}
  \label{tab:l04-prodcons-trace}
  \small
  \begin{tabular}{@{}cl>{\ttfamily}lccc@{}}
    \toprule
    Step & Event & \normalfont buffer & \texttt{add} & \texttt{rem} & \texttt{num} \\
    \midrule
     1 & P inserts 1                         & 1 - - -  & 1 & 0 & 1 \\
     2 & P inserts 2, 3, 4                   & 1 2 3 4  & 0 & 0 & 4 \\
     3 & P finds buffer full, waits on \texttt{c\_prod} & 1 2 3 4 & 0 & 0 & 4 \\
     4 & C removes 1, signals \texttt{c\_prod} & - 2 3 4 & 0 & 1 & 3 \\
     5 & P re-acquires \texttt{m}, re-tests, inserts 5 & 5 2 3 4 & 1 & 1 & 4 \\
     6 & C removes 2, 3, 4                   & 5 - - -  & 1 & 0 & 1 \\
     7 & C removes 5                         & - - - -  & 1 & 1 & 0 \\
     8 & C finds buffer empty, waits on \texttt{c\_cons} & - - - - & 1 & 1 & 0 \\
     9 & P inserts 6, signals \texttt{c\_cons} & - 6 - - & 2 & 1 & 1 \\
    10 & C re-acquires \texttt{m}, re-tests, removes 6 & - - - - & 2 & 2 & 0 \\
    \bottomrule
  \end{tabular}
\end{table}

\begin{keypoint}[Lesson summary]
  POSIX threads implement Birrell's mechanisms with concrete types and
  calls (\cref{tab:l04-birrell}).  \texttt{pthread\_create} starts a thread
  with one \texttt{void~*} argument and an attribute object that sets, for
  example, its stack size, scope and detach state; joinable threads must be
  joined, detached threads release their resources themselves, and
  \texttt{pthread\_exit} ends a thread without ending the process.
  Programs include \texttt{<pthread.h>}, are built with \texttt{-pthread}
  and check return values, and give every thread private argument data.
  Each shared resource has one globally visible mutex, locked in a global
  order and always unlocked; condition variables are waited on in a
  \texttt{while} loop and notified whenever a predicate may have changed,
  after unlocking where the decision to notify allows it, as in the
  producer-consumer program.
\end{keypoint}

\end{document}
